\section{Técnicas de Detección de Objetos}
Hasta el momento se han explorado las redes convolucionales como herramientas de clasificación:
dada una imagen de entrada, un modelo puede extraer características y asignar una etiqueta en
función de estas. Las redes convolucionales han sido ampliamente utilizadas durante décadas
en tareas como el reconocimiento de rostros y de peatones. Sin embargo, su aplicación directa
a la detección genérica de objetos es limitada.
\medskip

La detección genérica de objetos tiene como objetivo la localización y clasificación de múltiples
objetos en una imagen, etiquetándolos mediante una caja delimitadora (\textit{bounding box}) que
indica su ubicación. Este problema puede abordarse mediante redes
convolucionales utilizando un método de preprocesamiento conocido como ventana deslizante
(\textit{sliding window}). En este enfoque, se define una ventana de tamaño fijo que se desplaza
sobre la imagen, utilizando cada región cubierta como entrada independiente para la red.
\medskip

La principal desventaja de este método es su ineficiencia, ya que requiere evaluar múltiples
regiones de la imagen de forma redundante, lo que lo hace poco práctico para aplicaciones en
tiempo real. Además, el tamaño de los objetos detectables está limitado por el tamaño de la
ventana.
\medskip

Para abordar estas limitaciones, surgen dos principales enfoques para la detección de objetos.
El primero, conocido como enfoque de dos etapas (también denominado basado en regiones),
divide la tarea en una fase de propuesta de regiones y una fase de clasificación. El segundo,
conocido como enfoque de una etapa (también denominado basado en regresión o de tiempo real),
realiza de forma conjunta la clasificación y la predicción de las cajas delimitadoras a partir
de la imagen de entrada  \parencite{zhao2019object}.
\subsection{Detección de Dos Etapas}
\subsection{Detección de Una Etapa}